
@@M0@@, on ait @@M1@@), alors il en est de même pour @@M2@@.

Cette dernière remarque s'applique au cas où @@M3@@ est le quotient de @@M4@@ par un idéal régulier, car on peut alors prendre pour @@M5@@ le complexe de Koszul associé à un homomorphisme régulier définissant l'idéal de @@M6@@.

\subsection*{2.3.}
La structure de @@M7@@-algèbre définie en 2.1 sur @@M8@@ est fonctorielle en @@M9@@ et @@M10@@. On peut alors globaliser les constructions précédentes: soient @@M11@@ un schéma, @@M12@@, @@M13@@ deux @@M14@@-schémas; alors on peut localement sur @@M15@@ définir les morphismes (2.1.1), ils se recollent par fonctorialité, et définissent sur @@M16@@ une structure de @@M17@@-Algèbre graduée associative.

\textbf{Lemme 2.4.} \uline{Soient} @@M18@@ \uline{deux} @@M19@@-\uline{schémas}, \uline{et} @@M20@@ \uline{un} @@M21@@-\uline{Module plat}. \uline{Alors pour tout} @@M22@@, \uline{on a un isomorphisme canonique} :
@@M23@@

Il suffit de définir un isomorphisme fonctoriel dans le cas affine; on pourra alors par recollement le définir globalement dans le cas général.

Soient alors @@M24@@, @@M25@@, @@M26@@; @@M27@@ est défini par un @@M28@@-module plat @@M29@@. Soit @@M30@@ une résolution projective de @@M31@@ sur @@M32@@; les @@M33@@ sont les groupes d'homologie de @@M34@@. Comme @@M35@@ est plat sur @@M36@@, on a un isomorphisme canonique : @@M37@@. Par ailleurs, l'isomorphisme canonique : @@M38@@ donne un isomorphisme canonique : @@M39@@, d'où le résultat.

\textbf{Proposition 2.5.} \uline{Soient} @@M40@@ \uline{un schéma}, @@M41@@ \uline{et} @@M42@@ \uline{deux sous-schémas fermés de} @@M43@@, @@M44@@. \uline{On suppose que pour tout} @@M45@@, \uline{il existe un voisinage ouvert}
